\documentclass[10pt,a4paper]{amsart}
\usepackage[margin=19mm]{geometry}
\usepackage{amsmath,amssymb,mathtools}
\usepackage[scr=rsfs,cal=euler]{mathalfa}
\usepackage{xfrac}
\usepackage[unicode,hidelinks,bookmarks=false]{hyperref}
\newcommand{\ie}{i.e.\ }
\newcommand{\cf}{cf.\ }
\newcommand{\resp}{respectively\ }
\newcommand{\Subsection}{\subsection}
\providecommand{\nobreakdash}{\nobreak\hbox{-}}
\setlength{\emergencystretch}{2em}
\multlinegap0pt
\usepackage{amssymb}
\usepackage{xcolor}
\usepackage[shortlabels]{enumitem}
\newtheorem{lemma}{Lemma}
\newtheorem{claim}{Claim}
\title{\textbf{Social welfare optimisation under institutional reward and punishment}}
\author{Van An Nguyen, Vuong Khang Huynh, Huu Loi Bui, Hai Anh Ha, Quang Dung Le, Tan Dat Nguyen, Ngoc Ngu Nguyen, Zhao Song, Manh Hong Duong, Le Hong Trang, The Anh Han}
\date{Source: arXiv:2605.31330v1 (CC BY 4.0)}
\begin{document}
\maketitle

\noindent\emph{Source and licence.} \textbf{Social welfare optimisation under institutional reward and punishment}. Version arXiv:2605.31330v1 (CC BY 4.0), distributed by arXiv under the Creative Commons Attribution 4.0 International licence, http://creativecommons.org/licenses/by/4.0/, as recorded in the arXiv licence metadata for this exact version and shown on the abstract page https://arxiv.org/abs/2605.31330v1. Full licence notice: \texttt{licenses/arxiv-2605.31330-CC-BY-4.0.txt}.\par\smallskip

\noindent\textit{Editorial scope.} Counted unit: the article's lemma lem:lower\_bound\_p (source0.tex lines 1275--1279). The article's complete original proof of it is delivered with it (source0.tex lines 1280--1332, one \texttt{proof} environment of 53 lines). No mathematics is written, repaired or re-derived: the statement and the proof are the article's own lines, in the article's own notation, and no retained line is substituted or re-flowed.  Retained uncounted as the source-local context the proof needs: lines 988--990; lines 994--996.  Nothing was omitted from any retained span, so the package carries no omission record.  No retained line contains a citation, so the wrapper carries no bibliography and no bibliography entry.  No retained line contains a figure environment, an image-inclusion command or drawing code, so no image is delivered and the package declares no asset dependency.  Editorial formatting only, in addition: an AMS \texttt{amsart} preamble that restates the article's own declarations and macros used by the retained lines, namely the environments \texttt{lemma}, \texttt{claim} on separate counters, together with the standard packages those lines need.  The retained environments print in this document as Lemma 1, Claim 1 in order of appearance, whereas the article prints the same items with the numbering of its own full document.  The complete native source of the article as distributed in the arXiv e-print for this exact version is bundled unchanged as \texttt{sources/201713/source0.tex}.
\begin{lemma} \label{obs_eta0}
$S(u) > \eta_0$ for all $u > 0$.
\end{lemma}

\begin{lemma} \label{obs_lipschitz}
The function $S(u)$ is Lipschitz continuous and strictly increasing for all $u \in (0, 1)$.
\end{lemma}

\begin{lemma}[Left-side dominance by $\widehat{SW}(0)$]
\label{lem:lower_bound_p}

There exists a constant $p_l > 1$, independent of $\beta$, such that for all sufficiently large $\beta$, the following holds: If $\mu = \frac{-\delta}{\hat{a}} - \frac{\log p_l}{\hat{a}\beta}$, then $\widehat{SW}(\theta) \le \widehat{SW}(0)$ for all $\theta \in [0, \mu]$.
\end{lemma}

\begin{proof}
The proof is carried out similarly like in the reward case, with the only difference being that the quantity $(1-a)R(\theta)$ is replaced with $(1+\hat{a})\hat{R}(\theta)$.

% Similar to the reward case, we shall prove the following claims
% \begin{claim} \label{claim:decrease_near_zero_p}
% For sufficiently large $\beta$, $\widehat{SW}(\theta)$ is strictly decreasing on the interval $\hat{I}_0 := \left[0, \frac{1}{b\beta}\right]$.
% \end{claim}
% \begin{proof}
% We proceed by contradiction. Suppose there exists a point $\hat{\theta} \in I^*_0$ such that $\widehat{SW}'(\hat{\theta}) \ge 0$. Differentiating the objective yields:
% \[ 
% \widehat{SW}'(\theta) = \frac{N^2}{2}\left( (\delta + N\Delta)u'(\theta)S'(u(\theta)) - (1+b)\theta u'(\theta)\hat{S}'(u(\theta)) - (1+b)\hat{S}(u(\theta)) \right)
% \]

% For $\widehat{SW}'(\hat{\theta}) \ge 0$ to hold, we require:
% \[
% (\delta + N\Delta)u'(\theta)S'(u(\theta)) - (1+b)\theta u'(\theta)\hat{S}'(u(\theta)) \ge (1+b)\hat{S}(u(\theta))
% \]

% We know $u'(\hat{\theta}) = b\beta e^{\beta(b\hat{\theta} + \delta)} \le b\beta e^{\beta \delta + 1}$ for all $\hat{\theta} \in I_0$. Combining these facts with Lemma \ref{obs_eta0} and Lemma \ref{obs_lipschitz}, the inequality implies:
% \[
% (\delta + N\Delta)b\beta e^{\beta\delta + 1} L \geq (\delta + N\Delta)u'(\theta)S'(u(\theta)) - (1+b)\theta u'(\theta)\hat{S}'(u(\theta)) \ge (1+b) \hat{S}(u(\hat{\theta}))\geq (1+b)\eta_0.
% \]

% Since $\delta < 0$, we know that $\lim_{\beta \to \infty} \beta e^{\beta\delta} = 0$. Thus, for sufficiently large $\beta$, this inequality fails. Therefore, $\widehat{SW}'(\theta) < 0$, implying $\widehat{SW}(\theta) < \widehat{SW}(0)$ for all $\theta \in I^*_0$.
% \end{proof}

% \begin{claim} \label{claim:bound_away_from_zero_p}
% There exists an $\varepsilon_1 > 1$ independent of $\beta$, such that for all $\theta \in \left[\frac{1}{b\beta}, \frac{-\delta}{b} - \frac{\log \varepsilon_1}{b\beta}\right]$, $\widehat{SW}(\theta) \le \widehat{SW}(0)$.
% \end{claim}
% \begin{proof}
% First, note that this interval is valid for sufficiently large $\beta$. We require $\frac{1}{b\beta} \le \frac{-\delta}{b} - \frac{\log \varepsilon_1}{b\beta}$, which holds when $\beta \ge \frac{1 + \log \varepsilon_1}{-\delta}$.

% Suppose by contradiction that there exists a $\theta \in [0, \mu]$ such that $\widehat{SW}(\theta) > \widehat{SW}(0)$. This is equivalent to:
% \begin{align}
%     &(\delta + N\Delta) R(\theta) - (1+b)\theta\hat{S}(u(\theta)) > (\delta + N\Delta) R(0) \notag \\
%     \text{ iff } &(\delta + N\Delta)(R(\theta) - R(0)) > (1+b)\theta \hat{S}(u(\theta)) \label{ineq:lower_bound_p}
% \end{align}

% Since $R(\theta) - R(0) > 0$, applying Lemma \ref{obs_eta0} and Lemma \ref{obs_lipschitz} yields:
% \begin{align}
%     (1+b)\theta \eta_0 < (\delta + N\Delta)L(u(\theta) - u(0)) < (\delta + N\Delta)Lu(\theta) \label{ineq:3sides_p}
% \end{align}

% By making the substitution $\theta = \frac{-\delta}{b} - \frac{\log \varepsilon}{b\beta}$ (where $\varepsilon_1 \le \varepsilon \le e^{-\beta\delta - 1}$), we have $ u(\theta) = \varepsilon^{-1}$. Rearranging the left-most and right-most sides of the inequality, we have:
% \[
% \varepsilon \left(-\delta- \frac{\log \varepsilon}{\beta}\right) < \frac{(\delta + N\Delta)Lb}{(1+b)\eta_0}
% \]

% Observe that the left-hand side is an increasing function of $\varepsilon$, for all $1 \le \varepsilon \le e^{-\beta\delta - 1}$. Therefore, by choosing an $\varepsilon_1$ strictly independent of $\beta$ such that $-\delta \varepsilon_1 \ge \frac{(\delta + N\Delta)Lb}{(1+b)\eta_0}$, and a $\beta$ large enough such that $\varepsilon_1$ is valid (i.e. $1 \le \varepsilon_1 \le e^{-\beta\delta-1}$) the inequality fails for all valid $\varepsilon \ge \varepsilon_1$.
% \end{proof}

% Combining Claim \ref{claim:decrease_near_zero_p} and Claim \ref{claim:bound_away_from_zero_p}, setting $p_l = \varepsilon_1$ satisfies lemma \ref{lem:lower_bound_p}.
\end{proof}

\end{document}
